% ============================================================
% ch15.tex —— 第 15 章 瞬间的速度：发明导数
% ============================================================
\chapter{瞬间的速度：发明导数}

这一章处理一个深到骨头里的问题。它由伽利略四百年前用斜面实验逼问出来，答案由牛顿与莱布尼茨在十七世纪各自独立造出——而你将在本章亲手把答案\textbf{再发明一遍}。全程工具：初中代数，加第 14 章的比赛规则。不多一分。

\section{瞬间有速度吗？}

\begin{kaiwei}
骑车下坡，速度表指针一路上涨。在\textbf{某一瞬间}（比如第一秒末），你的速度是多少？
\end{kaiwei}

看似常识，细想头皮发麻。速度的定义是"路程 $\div$ 时间"。可"一瞬"没有时长：取瞬间 $= 0$ 秒，$0$ 秒里骑了 $0$ 米，速度 $= 0\div 0$——除以零，数学禁区。难道"瞬间的速度"像"方的圆"一样，是个天生的伪概念？

不是伪概念。速度表指针在每一刻都指着确定的读数，物理事实站在那里；出问题的是\textbf{我们的定义}——"路程除以时间"只够对付\textbf{匀速}运动，遇到变速就哑火。\textbf{旧定义遇到了做不了的事}——序章的发明套路启动：发一个新工具。

先看物理背景，问题不是凭空来的。伽利略在斜面实验（让铜球滚下精心打磨的斜槽，用水钟计时）中发现了一条惊人的规律：从静止下落的物体，走过的路程与时间的\textbf{平方}成正比：
\[
s = t^2 \qquad (\text{取合适单位；现实里近似 } s = \tfrac12 g t^2\text{，见本章末）}.
\]
验算：第一秒末走了 $1$ 米，第二秒末累计 $4$ 米，第三秒末累计 $9$ 米。各秒的\textbf{平均}速度：第 1 秒内 $1$ 米/秒，第 2 秒内 $3$ 米/秒，第 3 秒内 $5$ 米/秒——越来越快，而且每秒恰好多 $2$。奇妙的规律，但仍是"平均"的账。\textbf{在 $t$ 秒这一瞬间，速度表读数是多少？}只要能回答，"瞬间速度"就不是伪概念。带着问题发明新数学。

\section{先把平均速度算到极致}

手头唯一合法的工具是平均速度 $\dfrac{\text{路程差}}{\text{时间差}}$。那就用它，把时间间隔越切越短，盯住第一秒末（$t = 1$）：

\begin{center}
\begin{tabular}{ccc}
\toprule
时间区间 & 路程差 $s(\text{末}) - s(\text{初})$ & 平均速度 \\
\midrule
$[1, 2]$ & $4 - 1 = 3$ & $3 \div 1 = 3$ \\
$[1, 1.5]$ & $2.25 - 1 = 1.25$ & $1.25\div0.5 = 2.5$ \\
$[1, 1.2]$ & $1.44 - 1 = 0.44$ & $0.44 \div 0.2 = 2.2$ \\
$[1, 1.01]$ & $1.0201 - 1 = 0.0201$ & $0.0201\div0.01 = 2.01$ \\
$[1, 1.001]$ & $1.002001 - 1 = 0.002001$ & $\approx 2.001$ \\
$\cdots$ & $\cdots$ & $\to\, ?$ \\
\bottomrule
\end{tabular}
\end{center}

平均速度一路下滑：$3, 2.5, 2.2, 2.01, 2.001, \dots$——它\textbf{像是在逼近 $2$}，但每一个都\textbf{大于} $2$（区间从右侧切入，下坡越来越快，右侧的平均偏高）。从左侧也切一刀试试：区间 $[0.9, 1]$：$(1 - 0.81)\div0.1 = 1.9$；$[0.99, 1]$：$(1 - 0.9801)\div 0.01 = 1.99$——从下方逼近 $2$。两侧夹击，$2$ 米/秒呼之欲出。

但"呼之欲出"不是数学。硬数据只有有限行，谁能保证第 $10^{100}$ 行不会突然跳开？\textbf{一次性算到底}——别一行一行试，直接算\textbf{任意}小区间。设从 $t$ 秒到 $t+h$ 秒（$h$ 是"任你挑多小"的时间增量，可正可负）：
\[
\text{平均速度} = \frac{s(t+h) - s(t)}{h} = \frac{(t+h)^2 - t^2}{h} .
\]
展开分子——\textbf{初中乘法公式 $(a+b)^2 = a^2 + 2ab + b^2$ 正式上岗}，这是全书中学的知识发挥最大作用的一刻：
\[
(t+h)^2 - t^2 = \left(t^2 + 2th + h^2\right) - t^2 = 2th + h^2 = h\,(2t + h),
\]
于是（$h\ne0$ 时约去 $h$，完全合法）：
\[
\text{平均速度} = \frac{h(2t+h)}{h} = 2t + h \qquad (h \ne 0).
\]
现在把第 14 章的比赛开起来：对手随便挑 $h$（$0.1$？$0.001$？$10^{-100}$？），式子 $2t + h$ 的\textbf{落脚点是 $2t$}。结论钉死：

\[
\boxed{\;s = t^2 \text{ 时，} t \text{ 秒末的瞬间速度} = 2t\;}
\]

第一秒末 $2$ 米/秒（表格两侧夹击的正是它），第三秒末 $6$ 米/秒。\textbf{速度不再是孤立的读数，而是一个新的函数}：$v(t) = 2t$。路程函数 $s = t^2$ 经过这台机器，吐出了速度函数 $v = 2t$——一函数变另一函数，这就是"求导"的本质。

\begin{faming}
\textbf{土名}：瞬间变化率、瞬间速度 $\to$ \textbf{正式名}：\textbf{导数}（derivative），记号 $s'(t)$（拉格朗日记号，读"$s$ 一撇 $t$"）或 $\dfrac{ds}{dt}$（莱布尼茨记号，读"$s$ 对 $t$ 的导数"）。那台"平均变化率取落脚点"的机器叫\textbf{求导}。两条记号各自成家：$'$ 写起来快，$\dfrac{d}{dt}$ 在链式法则（第 16 章）里表意更清楚——本书随用随取。
\end{faming}

\section{几何意义：割线摇成切线}

把 $s = t^2$ 画在坐标系里（横轴 $t$、纵轴 $s$，一条抛物线）。平均速度 $\dfrac{s(t+h) - s(t)}{h}$ 是什么？——\textbf{连接两点 $(t, s(t))$ 与 $(t+h, s(t+h))$ 的直线的斜率}（纵坐标差除以横坐标差，初中一次函数斜率定义原样适用）。这条线叫\textbf{割线}（割过曲线两点的线）。

现在让 $h$ 进场比赛。$h$ 缩短，第二点沿曲线滑向第一点，割线绕着定点旋转——$h\to0$ 的刹那，割线摇身成为\textbf{切线}（只在该点贴住曲线的直线）。导数 $=$ 切线斜率：

\begin{tikzpicture}[scale=1.5]
  \draw[->] (-0.3,0) -- (2.6,0) node[right] {$t$};
  \draw[->] (0,-0.3) -- (0,4.8) node[above] {$s$};
  \draw[thick,domain=0:2.15,smooth] plot (\x,{\x*\x});
  \fill (1,1) circle (1.1pt);
  \draw[dashed,cheng!90!black,domain=0.12:2.05] plot (\x,{3*(\x-1)+1});
  \draw[dashed,molv!80!black,domain=0.05:1.85] plot (\x,{2.5*(\x-1)+1});
  \draw[very thick,shenlan,domain=0:1.85] plot (\x,{2*(\x-1)+1});
  \node[cheng!60!black] at (2.2,3.7) {\scriptsize $h=1$：割线，斜率 $3$};
  \node[molv] at (2.15,2.9) {\scriptsize $h=0.5$：割线，$2.5$};
  \node[shenlan] at (2.0,2.05) {\scriptsize $h\to0$：切线，$2$};
  \node at (0.5,3.3) {\scriptsize $s=t^2$};
\end{tikzpicture}

图上三条线与表格三行一一对应（斜率 $3, 2.5, 2$）。"切线"的初中定义是"与曲线只有一个公共点的直线"——这个定义其实有毛病（抛物线的切线确实只交一点，但正弦波的切线能交无穷多点；两条平行线……不切却只有一个"公共方向"）。\textbf{正确的定义恰恰是本书的路线：切线 $=$ 割线的极限位置}。又一次，"先玩后命名"造出了比旧定义更坚固的概念。

\section{直面 $0\div0$：我们没有作弊}

诚实的读者必须过这道坎：最后一步"约去 $h$ 再取落脚点 $2t$"，是不是偷偷"令 $h = 0$"的变体？\textbf{复盘每一分钟：}
\begin{enumerate}
\item 写出平均速度 $\dfrac{(t+h)^2 - t^2}{h}$：此刻 $h \ne 0$（否则除法无意义），合法；
\item 展开分子 $2th + h^2$，提取公因式 $h(2t+h)$：恒等变形，任何 $h$ 都合法；
\item 约去分子分母的 $h$：\textbf{前提 $h\ne0$}，合法——约完得到 $2t + h$，对一切 $h\ne0$ 成立；
\item 取 $h\to0$ 的落脚点：$h$ \textbf{从未等于零}，它只是"任你挑多小"；第 14 章的比赛规则担保 $2t + h$ 的归宿是 $2t$。
\end{enumerate}
全程没有除以零——\textbf{我们绕开了禁区，而不是冲了过去}。发明新数学，很多时候就是给"不能算的老路"修一条"能算的绕行道"。

顺带说历史：牛顿最初没把这件事说干净。他管 $h$ 叫"流量的瞬"，一会儿说它不是零（要除它），一会儿说它趋于零（要扔它）——被哲学家贝克莱主教抓住痛批："消失量的鬼魂！"（你凭什么对一会儿存在一会儿不存在的东西做除法？）这场论战逼得数学家们两百年后造出第 14 章的比赛规则（$\varepsilon$--$N$ 语言）。\textbf{你手里这台机器，是两百年争吵的最终成品}——而且你刚刚五分钟就操作正确了。

\section{机器量产：再喂三个函数}

这台机器当然不能只吃 $t^2$。同一套流程（展开、约分、取落脚点）再跑三台：

\paragraph{$s = t^3$。}分子 $(t+h)^3 - t^3$，用立方展开 $(t+h)^3 = t^3 + 3t^2h + 3th^2 + h^3$（可用乘法一步步验证）：
\[
\frac{3t^2h + 3th^2 + h^3}{h} = 3t^2 + 3th + h^2 \;\longrightarrow\; 3t^2 .
\]
\textbf{速度函数 $= 3t^2$。}（检验量纲直觉：三秒末速度 $27$ 米/秒，比 $t^2$ 的世界快得多——路程涨得快，速度也涨得快。）

\paragraph{$s = t$（匀速）。}分子 $(t+h) - t = h$，除以 $h$ 得 $1$——恒定。匀速的瞬间速度处处等于平均速度，废话，但机器认账（机器对匀速给出的答案与直觉一致，这是质检的重要一环）。

\paragraph{$s = \dfrac{1}{t}$。}
\[
\frac{\frac{1}{t+h} - \frac1t}{h} = \frac{\frac{t - (t+h)}{t(t+h)}}{h} = \frac{-h}{h\,t(t+h)} = \frac{-1}{t(t+h)} \;\longrightarrow\; -\frac{1}{t^2} .
\]
（通分、约分、取落脚点，三步全初中。）注意\textbf{负号}：函数在减小，导数是负的；$|v| = \tfrac1{t^2}$ 随 $t$ 增大而减小——减得越来越慢。\textbf{导数一条信息管两件事：变化的方向（正负）与快慢（大小）}。

数据汇总，规律探出鼻息：
\[
t^2 \to 2t, \qquad t^3 \to 3t^2, \qquad t^4 \to 4t^3\ (\text{习题}),\qquad t \to 1 = 1\cdot t^0 .
\]
大胆猜（下一章证明）：\textbf{$t^n$ 的导数是 $n\,t^{n-1}$（"指数落下来当系数，指数减一"）}。这条"幂法则"是求导流水线的第一件产品。

\section{回到物理：真正的自由落体}

现实中的自由落体是 $s = \tfrac12 g t^2$（$g \approx 9.8$ 米/秒$^2$，重力加速度）。求瞬间速度：分子展开 $\tfrac12 g\left[(t+h)^2 - t^2\right] = \tfrac12 g\,(2th + h^2)$，除以 $h$、取落脚点：
\[
v(t) = g\,t .
\]
伽利略从斜面实验里猜出的"速度与时间成正比"（他用水钟测的是平均速度，外推靠胆识），在这里\textbf{三行代数自动导出}——而且不止落体：任何"正比于 $t^2$"的运动（比如匀加速的汽车），速度必然线性增长。实验规律与代数推论在天上地下对齐，这是物理与数学四百年联盟的标准照。

\begin{cidian}
土名对照：瞬间变化率 $=$ \textbf{导数}；"求导机器"$=$ \textbf{求导/微分}；割线的极限位置 $=$ \textbf{切线}；平均变化率 $=$ \textbf{差商}（difference quotient，"差再除"）。导数是全微积分的第一主角，第 17 章的第二主角（积分）已在路上，两者将在第 18 章握手。
\end{cidian}

\begin{dongshou}
\begin{itemize}
  \yisi 用机器求 $s = t^2$ 在 $t = 3$ 处的瞬间速度，并造表验证：算区间 $[3, 3.5]$、$[3, 3.1]$、$[3, 3.01]$ 的平均速度（应依次为 $6.5, 6.1, 6.01$，从上方逼近 $6$）；再从左侧 $[2.9, 3]$、$[2.99, 3]$（$5.9, 5.99$）夹击。
  \yier 用 $(t+h)^4 = t^4 + 4t^3h + 6t^2h^2 + 4th^3 + h^4$ 求证 $s = t^4$ 的导数是 $4t^3$，确认幂法则第四票。（四次展开可由 $(t+h)^3(t+h)$ 乘出，动手乘一遍，感受"系数 $1, 4, 6, 4, 1$"的对称。）
  \yier 自由落体近似 $s = 5t^2$（米、秒，$g$ 取 $10$）。求速度函数 $v(t) = 10t$；从 $20$ 米高处落地需 $t = 2$ 秒（$5\times4 = 20$），落地前一刻速度 $= 20$ 米/秒。再算全程\textbf{平均}速度 $10$ 米/秒，与末速 $20$ 对比——匀加速运动的末速恰为平均速度的两倍，用 $v = gt$ 与 $s = \tfrac12 gt^2$ 解释为什么。
  \yier 求 $s = \sqrt t$ 在 $t = 4$ 处的瞬间速度。（提示：分子有理化——$\sqrt{t+h} - \sqrt t$ 乘共轭 $\sqrt{t+h}+\sqrt t$，得 $\dfrac{h}{\sqrt{t+h}+\sqrt t}$，除以 $h$ 后取落脚点 $\dfrac{1}{2\sqrt t}$；在 $t=4$ 处为 $\tfrac14$。幂法则对分数指数也灵：$(t^{1/2})' = \tfrac12 t^{-1/2}$。）
\end{itemize}
\end{dongshou}

\begin{shaonao}
\textbf{芝诺的第二次传唤："飞矢不动"。}芝诺说：飞行中的箭，在每一瞬间占据一个确定的位置——"和静止的箭一样"；每个瞬间都静止，全由瞬间组成的运动，岂能不是静止？运动是幻觉。

用本章的武器正面回击。芝诺偷换的概念：他默认"该瞬间的速度"由\textbf{该瞬间内的位移}决定（$0$ 秒 $0$ 米 $= 0\div0$，于是"静止"）。而本章造出的瞬间速度，定义是\textbf{以该瞬间为中心、两侧平均速度的共同落脚点}——它根本不是"那一瞬走了多远"，而是"那一瞬的趋势有多陡"。箭在飞行中的每个 $t$，$v(t) = gt > 0$：趋势处处非零，何来不动？

请你把这段反驳写成三句话的严格版，并指出它和"路标不是目的地"（第 13 章）、"极限是必胜的比赛"（第 14 章）共享哪一条哲学：\textbf{无穷过程的信息，在于结构（趋势、归宿），不在于逐点截取}。芝诺两千五百年前提出的这道题，答案的钥匙全书已经发给你三把。
\end{shaonao}

\begin{chengshi}
"落脚点存在"这件事我们对 $t^n$ 直接验收（展开式干净）；一般函数的导数可能不存在——图像有尖角（如 $|t|$ 在 $0$ 处，左右割线的落脚点一个 $-1$ 一个 $+1$，比赛打平不了）或断裂处没有切线，大学里这叫"可导性"，且"可导必连续、连续未必可导"。另外本章默认了"时间与路程可以无限细分"（物理上到普朗克尺度会成问题，但宏观世界足够真）——物理的诚实边界，与数学的理想机器，各管各的辖区。
\end{chengshi}

一台机器一台机器地造太慢。下一章，把作坊升级成流水线：\textbf{四条运算法则}，从此不用每次都展开 $(t+h)^n$。
